% Part: first-order-logic
% Chapter: syntax-and-semantics
% Section: terms-formulas

\documentclass[../../../include/open-logic-section]{subfiles}

\begin{document}

\olfileid{fol}{syn}{frm}

\olsection{الحدود وال\printtoken{P}{formula}}

متى أعطيت لغة رتبة أولى~$\Lang L$، أمكننا تعريف تعبيرات مبنية من
المعجم الأساسي لـ~$\Lang L$. ومن هذه التعبيرات، على وجه الخصوص،
\emph{الحدود} و\emph{!!{formula}s}.

\begin{defn}[الحدود]
\ollabel{defn:terms}
تُعرّف مجموعة \emph{الحدود}~$\Trm[L]$ في~$\Lang L$ استقرائيًا كما
يأتي:
\begin{enumerate}
\item كل !!{variable} حد.
\item كل !!{constant} في~$\Lang L$ حد.
\item إذا كانت $f$ !!{function} ذات $n$ من المحال، وكانت
  $t_1$، و\dots، و$t_n$ حدودًا، فإن
  $\Atom{f}{t_1, \ldots, t_n}$ حد.
\tagitem{limitClause}{
  لا شيء غير ذلك حد.}{}
\end{enumerate}
والحد الذي لا يحتوي على !!{variable}s \emph{حد مغلق}.
\end{defn}

\begin{explain}
تظهر ال!!{constant}s في تحديدنا للغة وللحدود بوصفها فئة مستقلة من
الرموز، لكن كان يمكن بدلًا من ذلك إدراجها بوصفها !!{function}s صفرية
المحل. وعندئذ أمكننا الاستغناء عن البند الثاني في تعريف الحدود. ولا
يلزمنا إلا أن نفهم $\Atom{f}{t_1, \ldots, t_n}$ على أنه $f$ وحده إذا
كان $n = 0$.
\end{explain}

\begin{defn}[الصيغ]
\ollabel{defn:formulas}
تُعرّف مجموعة \emph{!!{formula}s}~$\Frm[L]$ في اللغة~$\Lang L$
استقرائيًا على النحو الآتي:
\begin{enumerate}
\tagitem{prvFalse}{$\lfalse$ هي !!{formula} ذرية.}{}

\tagitem{prvTrue}{$\ltrue$ هي !!{formula} ذرية.}{}

\item إذا كان $R$ !!{predicate} ذا $n$ من المحال في~$\Lang L$، وكانت
  $t_1$، و\dots، و$t_n$ حدودًا في~$\Lang L$، فإن
  $\Atom{R}{t_1,\ldots, t_n}$ !!{formula} ذرية.

\item إذا كان $t_1$ و$t_2$ حدين في~$\Lang L$، فإن
  $\Atom{\eq}{t_1, t_2}$ !!{formula} ذرية.

\tagitem{prvNot}{إذا كانت $!A$ !!a{formula}، فإن $\lnot !A$
  !!a{formula}.}{}

\tagitem{prvAnd}{إذا كانت $!A$ و$!B$ !!{formula}s، فإن $(!A \land
  !B)$ !!a{formula}.}{}

\tagitem{prvOr}{إذا كانت $!A$ و$!B$ !!{formula}s، فإن $(!A \lor !B)$
  !!a{formula}.}{}

\tagitem{prvIf}{إذا كانت $!A$ و$!B$ !!{formula}s، فإن $(!A \lif !B)$
  !!a{formula}.}{}

\tagitem{prvIff}{إذا كانت $!A$ و$!B$ !!{formula}s، فإن
  $(!A \liff !B)$ !!a{formula}.}{}

\tagitem{prvAll}{إذا كانت $!A$ !!a{formula} وكان $x$ !!a{variable}،
  فإن $\lforall[x][!A]$ !!a{formula}.}{}

\tagitem{prvEx}{إذا كانت $!A$ !!a{formula} وكان $x$ !!a{variable}،
  فإن $\lexists[x][!A]$ !!a{formula}.}{}

\tagitem{limitClause}{لا شيء غير ذلك !!a{formula}.}{}
\end{enumerate}
\end{defn}

\begin{explain}
تعريف مجموعة الحدود وتعريف مجموعة ال!!{formula}s
\emph{تعريفان استقرائيان}. ففي جوهر الأمر، نبني مجموعة
ال!!{formula}s عبر عدد لا نهائي من المراحل. وفي المرحلة الابتدائية
نعلن أن جميع الصيغ الذرية صيغ؛ وهذا يقابل الحالات الأولى القليلة من
التعريف، أي حالات
\iftag{prvTrue}{$\ltrue$، }{}%
\iftag{prvFalse}{$\lfalse$، }{}%
$\Atom{R}{t_1,\dots,t_n}$ و$\Atom{\eq}{t_1,t_2}$. وهكذا تعني
«!!{formula} ذرية» أي !!{formula} من هذا الشكل.

وتعطي الحالات الأخرى من التعريف قواعد لبناء !!{formula}s جديدة من
!!{formula}s سبق بناؤها. وفي المرحلة الثانية يمكننا استعمال تلك
القواعد لبناء !!{formula}s من !!{formula}s ذرية. وفي المرحلة الثالثة
نبني صيغًا جديدة من الصيغ الذرية ومن الصيغ التي حصلنا عليها في
المرحلة الثانية، وهكذا. ويُعد الشيء !!{formula} إذا بُني في نهاية
المطاف في مرحلة من هذه المراحل، ولا يُعد أي شيء آخر كذلك.
\end{explain}

نكتب، اصطلاحًا، $\eq$ بين حجتيها ونحذف القوسين:
فـ~$\eq[t_1][t_2]$ اختصار لـ~$\Atom{\eq}{t_1,t_2}$. ويُختصر أيضًا
$\lnot \Atom{\eq}{t_1,t_2}$ إلى~$\eq/[t_1][t_2]$. وعند كتابة صيغة
$(!B \ast !C)$ منشأة من $!B$ و$!C$ باستعمال رابط ثنائي
المحل~$\ast$، سنحذف غالبًا الزوج الخارجي من الأقواس ونكتب
ببساطة~$!B \ast !C$.

\begin{intro}
تشترط بعض كتب المنطق أن يرد ال!!{variable}~$x$ في~$!A$ لكي تدخل
\iftag{prvEx}{$\lexists[x][!A]$ }{}%
\iftag{notprvEx,notprvAll}{}{و }%
\iftag{prvAll}{$\lforall[x][!A]$ }{}%
في عداد
\iftag{notprvEx,notprvAll}{!!a{formula}}{!!{formula}s}.
ولا يحدث ما يضر إذا لم تشترط ذلك، كما أن تركه ييسر الأمور.
\end{intro}

\begin{tagblock}{defNot,defOr,defAnd,defIf,defIff,defTrue,defFalse,defEx,defAll}
\begin{defn}
يجب فهم الصيغ المنشأة باستعمال المؤثرات المعرّفة على النحو الآتي:

\begin{tagenumerate}{defTrue,defFalse,defNot,defOr,defAnd,defIf,defIff,defEx,defAll}

\tagitem{defTrue}{$\ltrue$ اختصار لـ
  \iftag{prvFalse}{$\lnot\lfalse$}{$(!A \lor \lnot !A)$، حيث تكون
    !!{formula}~$!A$ ذرية ثابتة ما}.}{}

\tagitem{defFalse}{$\lfalse$ اختصار لـ
  \iftag{prvTrue}{$\lnot\ltrue$}{$(!A \land \lnot !A)$، حيث تكون
    !!{formula}~$!A$ ذرية ثابتة ما}.}{}

\tagitem{defNot}{$\lnot !A$ اختصار لـ~$!A \lif \lfalse$.}{}

\tagitem{defOr}{$!A \lor !B$ اختصار لـ
  \iftag{prvAnd}{$\lnot(\lnot !A \land \lnot !B)$}{$\lnot !A \lif
    !B$}.}{}

\tagitem{defAnd}{$!A \land !B$ اختصار لـ
  \iftag{prvOr}{$\lnot(\lnot !A \lor \lnot !B)$}{$\lnot (!A \lif
    \lnot !B)$}.}{}

\tagitem{defIf}{$!A \lif !B$ اختصار لـ
  \iftag{prvOr}{$(\lnot !A \lor !B)$}{$\lnot (!A \land \lnot !B)$}.}{}

\tagitem{defIff}{$!A \liff !B$ اختصار لـ~$(!A \lif !B) \land (!B
  \lif !A)$.}{}

\tagitem{defAll}{$\lforall[x][!A]$ اختصار لـ~$\lnot\lexists[x][\lnot !A]$.}{}

\tagitem{defEx}{$\lexists[x][!A]$ اختصار لـ~$\lnot\lforall[x][\lnot !A]$.}{}
\end{tagenumerate}
\end{defn}
\end{tagblock}

إذا عملنا في لغة لتطبيق معين، كتبنا كثيرًا ال!!{predicate}s
وال!!{function}s ثنائية المحل بين الحدين المعنيين؛ مثل
$t_1 < t_2$ و$(t_1 + t_2)$ في لغة الحساب، و$t_1 \in t_2$ في لغة
نظرية المجموعات. بل من الاصطلاح كتابة دالة الخلف في لغة الحساب
\emph{بعد} حجتها:~$t'$. لكن هذه رسميًا ليست إلا اختصارات اصطلاحية،
على الترتيب، لـ~$\Atom{\Obj A^2_0}{t_1, t_2}$ و
$\Obj f^2_0(t_1, t_2)$ و$\Atom{\Obj A^2_0}{t_1, t_2}$ و$\Obj f^1_0(t)$.

\begin{defn}[الهوية التركيبية]
يعبّر الرمز~$\ident$ عن الهوية التركيبية بين سلاسل الرموز؛ أي إن
$!A \ident !B$ إذا وفقط إذا كان $!A$ و$!B$ سلسلتين من الرموز لهما
الطول نفسه وتحتويان الرمز نفسه في كل موضع.
\end{defn}

يمكن أن تحف بالرمز~$\ident$ سلاسل ناتجة من الوصل؛ فمثلًا تعني
$!A \ident (!B \lor !C)$ أن سلسلة الرموز~$!A$ هي السلسلة نفسها الناتجة
من وصل قوس فاتح، والسلسلة~$!B$، والرمز~$\lor$، والسلسلة~$!C$، وقوس
غالق، بهذا الترتيب. وإذا كان الأمر كذلك، علمنا أن الرمز الأول من~$!A$
قوس فاتح، وأن~$!A$ تحتوي على~$!B$ بوصفها سلسلةً جزئية (تبدأ عند الرمز
الثاني)، وأن هذه السلسلة الجزئية يليها~$\lor$، وهكذا.

لما كانت الحدود وال!!{formula}s تُبنى من عناصر أساسية بواسطة تعريفات
استقرائية، أمكننا استعمال مبدأي الاستقراء الآتيين لإثبات قضايا عنها.

\begin{lem}[\emph{مبدأ الاستقراء على الحدود}]
    \ollabel{lem:trmind}
    لتكن $\Lang L$ لغة رتبة أولى.
    إذا كانت خاصية ما~$P$ بحيث
    %
    \begin{enumerate}
        \item تصدق في كل !!{variable}~$v$،
        %
        \item وتصدق في كل !!{constant}~$a$ من~$\Lang L$،
        %
        \item وتصدق في $f(t_1,\dotsc,t_n)$ متى صدقت في
        $t_1$، و\dots، و$t_n$ وكانت $f$ !!{function} ذات $n$ من
            المحال في~$\Lang L$،
    \end{enumerate}
    (على افتراض أن $t_1$، و\dots، و$t_n$ حدود في~$\Lang{L}$)،
    فإن $P$ تصدق في كل حد في~$\Trm[L]$.
\end{lem}

\begin{prob}
    أثبت \olref[fol][syn][frm]{lem:trmind}.
\end{prob}

\begin{lem}[\emph{مبدأ الاستقراء على ال!!{formula}s}]
    \ollabel{thm:frmind}
    لتكن $\Lang L$ لغة رتبة أولى.
    إذا كانت خاصية ما~$P$ تصدق في جميع ال!!{formula}s الذرية
    وكانت بحيث
    %
    \begin{enumerate}
        \tagitem{prvNot}{تصدق في $\lnot !A$ متى صدقت
            في~$!A$؛}{}
        \tagitem{prvAnd}{تصدق في $(!A \land !B)$
            متى صدقت في $!A$ و~$!B$؛}{}
        \tagitem{prvOr}{تصدق في $(!A \lor !B)$
            متى صدقت في $!A$ و~$!B$؛}{}
        \tagitem{prvIf}{تصدق في $(!A \lif !B)$
            متى صدقت في $!A$ و~$!B$؛}{}
        \tagitem{prvIff}{تصدق في $(!A \liff !B)$
            متى صدقت في $!A$ و~$!B$؛}{}
        \tagitem{prvEx}{تصدق في $\lexists[x][!A]$
            متى صدقت في~$!A$؛}{}
        \tagitem{prvAll}{تصدق في $\lforall[x][!A]$
            متى صدقت في~$!A$؛}{}
  \end{enumerate}
  (على افتراض أن $!A$ و$!B$ !!{formula}s في~$\Lang{L}$)،
  فإن $P$ تصدق في جميع الصيغ في~$\Frm[L]$.
\end{lem}

\end{document}
